\section{Migrado de capitulo7\_implementacao\_do\_algoritmo\_genetico.tex}

A implementação do algoritmo genético Sunny no sistema IGRIS representa a aplicação prática dos conceitos apresentados no capítulo anterior. Enquanto o Capítulo 3 discutiu os fundamentos dos algoritmos genéticos e a forma como princípios inspirados na evolução biológica podem ser transpostos para a computação, este capítulo concentra-se em explicar como esses conceitos foram utilizados, de maneira concreta, na geração de grades horárias acadêmicas.

O problema enfrentado pelos algoritmos está inserido no contexto do \textit{University Course Timetabling Problem} (UCTP), que envolve a distribuição de demandas acadêmicas em horários disponíveis, considerando restrições como turma, turno, professor habilitado, disponibilidade docente e conflitos de alocação. Trata-se de um problema de busca e otimização, pois o número de combinações possíveis aumenta significativamente conforme crescem a quantidade de turmas, disciplinas, professores e horários disponíveis. Diante desse cenário, uma busca exaustiva torna-se pouco viável, o que justifica o uso de métodos aproximativos, como os algoritmos genéticos.

Como discutido no Capítulo 3, algoritmos genéticos não têm como finalidade garantir, necessariamente, a obtenção da melhor solução possível em termos absolutos. Seu objetivo é encontrar soluções satisfatórias dentro de um tempo computacional aceitável. Essa característica está presente nos algoritmos do IGRIS, que trabalham com sucessivas gerações de grades candidatas, avaliando sua qualidade por meio de uma função \textit{fitness} e preservando, ao longo da execução, as melhores soluções encontradas. Assim, a resposta final do algoritmo corresponde à melhor grade identificada de acordo com os critérios definidos pela implementação.

\section{Implementação do Algoritmo Sunny}
\label{sec:7.1}

No Sunny, a execução do algoritmo genético ocorre no serviço responsável pela geração da grade horária. Esse serviço recebe uma requisição contendo os dados acadêmicos necessários ao processo de alocação, como turmas, matrizes curriculares, disponibilidades docentes e vínculos entre usuários e disciplinas. A partir dessas informações, o algoritmo constrói uma estrutura de contexto denominada EvolutionContext, utilizada para organizar os parâmetros, os mapas auxiliares e os caches necessários durante uma execução específica.

Essa organização é importante porque o algoritmo genético realiza avaliações repetidas. A cada geração, vários indivíduos são avaliados, selecionados, recombinados e modificados. Caso todas essas verificações fossem feitas diretamente sobre as listas originais da requisição, o custo computacional tenderia a ser maior e a implementação ficaria menos eficiente. Por esse motivo, o Sunny cria estruturas de acesso rápido, como mapas de disponibilidades, matrizes, turmas, professores habilitados por disciplina, disponibilidades por turno, turno obrigatório de cada matriz, usuário associado a cada disponibilidade e dia da semana correspondente a cada horário.

Essa etapa pode ser compreendida como uma adaptação prática do conceito de ambiente apresentado no Capítulo 3. Na biologia, o ambiente é o espaço em que os organismos competem, sobrevivem e são selecionados. No algoritmo genético, esse ambiente é representado pelo próprio problema e por suas restrições. No caso do Sunny, ele é formado pelas regras de alocação acadêmica: uma matriz curricular precisa ser associada a uma disponibilidade compatível, o professor precisa estar habilitado para a disciplina, o turno da disponibilidade deve corresponder ao turno da turma e conflitos de horário precisam ser evitados.

A implementação estabelece que os valores utilizados pelo microsserviço sejam definidos por meio de parâmetros armazenados na tabela de configuração do sistema principal, permitindo ao administrador personalizar o comportamento do algoritmo conforme as necessidades institucionais. Atualmente, o banco de dados possui valores padrão definidos a partir de um extenso processo de testes de desempenho, os quais serão detalhados em seção específica, em que diferentes configurações foram avaliadas com o objetivo de identificar aquelas que apresentavam melhor performance, menor incidência de erros, menor tempo de execução e menor consumo de recursos durante a geração das grades.

A forma como uma solução é representada é uma das decisões mais relevantes em um algoritmo genético. No Sunny, cada gene é composto por dois elementos principais: matrizId e disponibilidadeId. O campo matrizId identifica a demanda acadêmica que precisa ser alocada, enquanto o campo disponibilidadeId indica a disponibilidade escolhida para essa demanda. O cromossomo, por sua vez, é formado por uma lista desses genes e representa uma grade horária candidata completa. Essa codificação traduz o problema real para uma estrutura que pode ser manipulada computacionalmente. No domínio acadêmico, o sistema precisa montar uma grade de aulas. No domínio genético, o algoritmo trabalha com listas de pares entre matrizes e disponibilidades. Quando uma demanda não pode ser alocada de maneira compatível, o Sunny utiliza o valor ID\_NAO\_ALOCADO = 0.

A população inicial do Sunny é formada por cromossomos gerados a partir dos dados recebidos na requisição. Essa geração não ocorre de forma totalmente aleatória. Para cada matriz curricular, o Sunny verifica inicialmente o turno obrigatório associado à turma, busca as disponibilidades compatíveis com esse turno, identifica os professores habilitados para a disciplina correspondente e filtra as disponibilidades pertencentes a esses professores. Somente após essa etapa de filtragem o algoritmo escolhe aleatoriamente uma disponibilidade entre as alternativas consideradas válidas. Essa abordagem pode ser classificada como uma inicialização aleatória guiada por restrições, aproximando-se da ideia de codificação restritiva discutida por \citeonline{michalewicz1996}.

A função \textit{fitness} é o mecanismo utilizado pelo Sunny para medir a qualidade de cada cromossomo. No Sunny, a avaliação começa com uma pontuação base proporcional ao número de genes do cromossomo. Em seguida, a função percorre os genes e aplica penalidades ou bônus conforme as condições encontradas. Quando um gene está marcado como não alocado, o indivíduo recebe uma penalidade padrão elevada. O Sunny também verifica conflitos de professor por meio de uma chave composta por usuário, dia e turno. Se mais de uma matriz for associada ao mesmo professor no mesmo período, o algoritmo identifica uma colisão e aplica uma penalidade proporcional à quantidade de conflitos. A função \textit{fitness} também aplica um bônus de diversidade associado à quantidade de matrizes atribuídas a um professor, com valor padrão significativamente menor que a penalidade por conflito, de modo que ele atua como critério secundário sem se sobrepor à prioridade de reduzir alocações inválidas.

Após a avaliação, o Sunny ordena os indivíduos e aplica duas estratégias: elitismo e seleção por torneio. O elitismo preserva diretamente uma parcela dos melhores indivíduos para a geração seguinte, evitando que boas soluções sejam perdidas por alterações dos operadores evolutivos. A seleção dos pais é realizada por torneio, no qual o algoritmo escolhe aleatoriamente um grupo de indivíduos e seleciona o melhor entre eles, criando pressão seletiva com grau de aleatoriedade que preserva a diversidade populacional.

O \textit{crossover} do Sunny utiliza uma estratégia próxima do \textit{crossover} uniforme, em que o método percorre cada posição do cromossomo e decide aleatoriamente se o gene daquela posição será herdado do primeiro ou do segundo pai. A mutação é aplicada gene a gene: para cada gene, o algoritmo sorteia uma probabilidade e, caso seja inferior à taxa configurada, substitui a disponibilidade associada àquela matriz. Quando a mutação ocorre, o Sunny não escolhe uma nova disponibilidade de forma irrestrita; o algoritmo consulta novamente o turno obrigatório e filtra os professores habilitados antes de escolher aleatoriamente uma alternativa válida.

Após a recombinação e a mutação, o Sunny executa um mecanismo de reparo parcial que identifica situações em que uma mesma turma possui mais de uma alocação no mesmo dia e turno. Quando isso ocorre, o algoritmo preserva a primeira ocorrência e marca as demais como não alocadas. O Sunny não realiza reparo por realocação; o que existe é um reparo parcial por invalidação de alocações conflitantes. Essa abordagem combina duas estratégias de tratamento de restrições: reparo e penalização.

Ao final da execução, o melhor cromossomo encontrado é decodificado em uma lista de alocações de grade. Na decodificação, cada gene é transformado em uma estrutura de alocação contendo turma, matriz, disciplina, disponibilidade, professor e status. O Sunny calcula então as métricas finais, que contabilizam disciplinas não alocadas, inconsistências de turno e conflitos de horário de professor. O ciclo evolutivo é executado até atingir o número máximo de gerações definido nos parâmetros, registrando ao longo da execução a geração em que uma solução sem erros é encontrada, sem que isso interrompa o laço evolutivo.

\section{Implementação do Algoritmo Morpheus}
\label{sec:7.2}

No Morpheus, a execução do algoritmo genético ocorre principalmente na classe GradeService, responsável por receber os dados acadêmicos, organizar as estruturas internas de processamento, gerar a população inicial, aplicar os operadores evolutivos e retornar a melhor grade horária encontrada. A entrada principal do algoritmo é uma requisição do tipo GradeRequest, composta por disponibilidades docentes, matrizes curriculares, turmas, vínculos entre professores e disciplinas, parâmetros do algoritmo genético e parâmetros da função \textit{fitness}.

A obtenção dos dados pode ocorrer por meio de uma camada externa ao algoritmo, responsável por consultar o banco de dados e montar a estrutura de entrada. No arquivo data\_service.py, são buscadas disponibilidades, matrizes, turmas e vínculos entre usuários e disciplinas. Entretanto, a execução do algoritmo genético em si não realiza consultas diretas ao banco: o GradeService recebe uma requisição já estruturada, constrói seus índices e caches em memória e, a partir dessas informações, executa o processo evolutivo.

Antes de iniciar a evolução, o Morpheus constrói índices em memória por meio do método \_construir\_indices. Essa etapa transforma as listas recebidas na requisição em estruturas de acesso rápido, criando mapas para recuperar disponibilidades, matrizes e turmas por seus identificadores, além de caches que relacionam disciplinas aos professores habilitados, turnos às disponibilidades e associações entre matrizes, turmas e turnos. Entre essas estruturas, destaca-se o cache de candidatos por matriz: para cada matriz curricular, o algoritmo identifica quais disponibilidades podem ser utilizadas, considerando o turno da turma e o vínculo entre professor e disciplina.

A representação genética do Morpheus utiliza explicitamente uma dataclass chamada Gene, com dois campos: matriz\_id e disponibilidade\_id. O cromossomo, formado por uma lista de genes, corresponde a uma grade horária candidata completa. Além da classe Gene, o Morpheus define a dataclass Individuo, que armazena o cromossomo e sua respectiva \textit{fitness}, e a dataclass Metricas, que concentra informações como quantidade de erros, conflitos, disciplinas não alocadas e dados inválidos. Quando uma matriz não possui candidatos válidos, o Morpheus utiliza o valor ID\_NAO\_ALOCADO = 0.

A população inicial é gerada pelo método \_inicializar\_populacao, que cria cromossomos por meio do método \_gerar\_cromossomo. Para cada matriz, o algoritmo verifica se existem turno, disciplina e turma associados e consulta a lista de candidatos válidos previamente construída. Caso existam candidatos, o algoritmo escolhe aleatoriamente uma disponibilidade, tentando priorizar aquelas que ainda não geram colisão para a mesma turma no mesmo dia e turno. Essa estratégia pode ser classificada como uma inicialização aleatória guiada por restrições.

A função \textit{fitness} é implementada no método \_calcular\_fitness. A avaliação começa com uma pontuação base proporcional ao tamanho do cromossomo, acrescida de bônus por alocações válidas e bônus de diversidade. As penalidades são os elementos mais decisivos: genes não alocados recebem penalidade padrão de 10.000 pontos; conflitos de professor recebem 12.000 pontos, por serem considerados os mais críticos; conflitos de turma recebem 8.000 pontos; dados inválidos recebem 10.000 pontos. Essa hierarquia de pesos demonstra uma priorização interna da implementação.

A evolução da população ocorre por meio do método \_evoluir\_populacao, que combina elitismo, seleção por torneio com \textit{crossover} de um ponto, mutação gene a gene, melhoria local e reparo inteligente. O \textit{crossover} de um ponto sorteia uma posição de corte no cromossomo, de modo que o primeiro filho recebe os genes iniciais do primeiro pai e os genes finais do segundo. A mutação é guiada por restrições do domínio, selecionando a nova disponibilidade dentro do conjunto de candidatos válidos da matriz correspondente.

Uma característica própria do Morpheus é a existência de uma etapa explícita de melhoria local antes do reparo inteligente. Após a mutação, o método \_melhorar\_alocacoes\_locais percorre o cromossomo buscando genes não alocados e tenta encontrar um candidato capaz de preencher a lacuna com pontuação interna de conflito igual a zero. Em seguida, o método \_reparar\_inteligente percorre o cromossomo em até duas passagens, identificando genes problemáticos e buscando o melhor candidato dentro das disponibilidades válidas por meio de uma pontuação interna de conflito. Quando encontra uma opção sem conflito, realiza a substituição; quando não há solução perfeita, pode aceitar uma melhoria parcial.

O critério de parada do Morpheus combina um limite fixo de gerações com mecanismos de parada antecipada. A execução pode ser interrompida quando uma solução sem erros é encontrada, ou quando o algoritmo passa muitas gerações sem melhorar a melhor \textit{fitness} global. O limite de gerações sem melhoria é calculado como 20\% do total configurado, respeitando um mínimo de 25 e um máximo de 200 gerações. Ao final, o melhor cromossomo é decodificado pelo método \_decodificar\_cromossomo e as métricas finais são calculadas pelo método \_calcular\_metricas\_detalhadas.

\subsection{Implementação do Algoritmo Rocket}
\label{sec:7.3}

No Rocket, a execução do algoritmo genético ocorre a partir de uma requisição que concentra os dados acadêmicos necessários à geração da grade horária, contendo informações sobre disponibilidades docentes, matrizes curriculares, turmas, vínculos entre professores e disciplinas, parâmetros do algoritmo genético, parâmetros da função \textit{fitness} e, quando necessário, uma URL de callback para envio do resultado.

Antes de iniciar a evolução, o Rocket organiza os dados recebidos por meio da construção de índices auxiliares, criando mapas de acesso rápido para disponibilidades, matrizes, turmas, professores, dias da semana, turnos e vínculos. Entre essas estruturas, destaca-se a lista de candidatos, que armazena para cada matriz curricular quais disponibilidades podem ser utilizadas, considerando o turno da turma e o vínculo entre professor e disciplina.

A representação genética do Rocket difere da utilizada no Sunny e no Morpheus. Enquanto nesses o gene é representado por um par explícito entre matriz e disponibilidade, no Rocket o gene é armazenado como um único valor inteiro correspondente ao identificador de uma disponibilidade. A posição desse gene dentro do cromossomo possui significado próprio: a implementação mantém uma lista ordenada de matrizes, e cada posição do vetor de genes corresponde a uma matriz específica. O valor presente naquela posição indica qual disponibilidade foi atribuída à matriz correspondente. Genes sem alocação recebem o valor unallocatedID = 0.

A população inicial é gerada a partir dos candidatos previamente definidos para cada matriz. Essa geração pode ser classificada como uma inicialização híbrida: a escolha final é aleatória, mas ocorre preferencialmente dentro de um subconjunto menos propenso a conflitos. Quando há alternativas livres para a turma, o sorteio ocorre entre essas opções; quando não há, entre todos os candidatos disponíveis. Um aspecto técnico importante é que a geração da população inicial pode ser paralelizada por meio de \textit{goroutines}, melhorando o desempenho operacional em execuções com populações maiores.

A função \textit{fitness} do Rocket calcula uma pontuação inicial a partir do número total de genes multiplicado pela pontuação base. Além disso, cada alocação válida recebe um bônus adicional (validBonus) e é aplicado um bônus de diversidade (diversityBonus) de papel secundário em relação às penalidades. As penalidades são os elementos mais decisivos: conflito de professor possui o maior peso, refletindo que um docente não pode ministrar duas aulas simultaneamente; conflito de turma possui peso menor; genes não alocados e dados inválidos também recebem penalidades elevadas. Os pesos são configuráveis por meio dos parâmetros da requisição.

A seleção por torneio cria pressão seletiva moderada. O \textit{crossover} utilizado é de um ponto, implementado de modo que dois pais gerem um novo filho: o algoritmo sorteia um ponto de corte e os genes anteriores a esse ponto são herdados do primeiro pai, enquanto os posteriores são herdados do segundo. A mutação é aplicada gene a gene, selecionando a nova disponibilidade dentro da lista de candidatos válidos para a matriz correspondente.

O reparo do Rocket é implementado como uma tentativa ativa de realocação. O algoritmo calcula ocupações de professor e turma, identifica genes associados a conflitos ou ausência de alocação, e busca a melhor alternativa possível dentro do conjunto de candidatos. A escolha da nova disponibilidade é orientada por uma pontuação interna de conflito: candidatos que colidem com a agenda do professor recebem acréscimo maior, enquanto candidatos que colidem com a agenda da turma recebem acréscimo menor. O algoritmo escolhe o candidato com menor pontuação.

O critério de parada do Rocket combina limite fixo de gerações com duas formas de parada antecipada: quando uma solução sem erros é encontrada, e quando o algoritmo passa muitas gerações sem melhorar a melhor \textit{fitness} global. O limite de gerações sem melhoria é calculado com base no total configurado, respeitando valores mínimo e máximo. Ao final, o melhor cromossomo é decodificado em uma grade horária e as métricas finais são calculadas, contabilizando disciplinas não alocadas, conflitos de professor, conflitos de turma e dados incompletos, gerando também mensagens textuais de conflito para integração com outros sistemas.